\section{A larger exact class: quadratic stars with affine rows}\label{sec:star}

The obstruction in Section~\ref{sec:overlap} motivates support directions
that span adjacent pairs. A tractable extension is a center variable $y$
with leaves $x_1,\ldots,x_k$. The objective is
\begin{equation}\label{eq:star}
q(y,x)=a_0+a_1y+a_2y^2+
\sum_{i=1}^k\{d_ix_i^2+(e_iy+f_i)x_i\}.
\end{equation}
All objective coefficients, affine-row coefficients, right-hand sides,
and finite variable bounds are rational; finite binary floating-point
inputs mean their exact rational values. Each supplied affine row contains
only $y$ and at most one leaf. Rows containing only $y$ are allowed;
rows involving two different leaves are outside this class. Every
quadratic coefficient can have either sign. The same structure must hold
for the scalar support objective obtained from a proposed direction.

For fixed $y$, each leaf is constrained to an interval
\begin{equation}\label{eq:leaf-bounds}
L_i(y)\le x_i\le U_i(y),\qquad
L_i(y)=\max_{r\in\mathcal L_i}(\alpha_{ir}+\beta_{ir}y),\quad
U_i(y)=\min_{r\in\mathcal U_i}(\gamma_{ir}+\delta_{ir}y).
\end{equation}
The original bounds provide at least one function in each envelope.
Projecting each bounded center--leaf polygon onto the center gives a
closed rational interval. Their intersection, together with the center
bounds and center-only rows, is exactly the feasible center interval
$I$: conditional leaf feasibility is independent across leaves. An
empty intersection proves this support domain empty.

\begin{theorem}[Exact rational star support]\label{thm:star}
For the nonempty domain just described, the minimum of
\eqref{eq:star} is rational and has a rational minimizer. It is obtained
by a finite rational partition of $I$ on which each leaf has an affine
minimizer and the conditional optimum is quadratic in $y$. Both the
partition and the attained minimum can be reconstructed exactly from
the supplied rows, bounds, and coefficients.
\end{theorem}
\begin{proof}
Partition $I$ at all intersections of affine functions in the envelopes
of each leaf. On every open resulting interval, the active functions
$L_i$ and $U_i$ are fixed and affine.

If $d_i>0$, the unconstrained leaf minimizer is the affine function
$v_i(y)=-(e_iy+f_i)/(2d_i)$. The constrained minimizer is
$\max\{L_i(y),\min\{v_i(y),U_i(y)\}\}$. Further partition at
intersections of $v_i$ with its two active bounds. Each remaining
interval has one affine minimizing rule.

If $d_i\le0$, a minimizing leaf value is an endpoint. The difference
between the upper and lower endpoint objectives factors as
\[
q_i(y,U_i(y))-q_i(y,L_i(y))
=(U_i(y)-L_i(y))
\{d_i(U_i(y)+L_i(y))+e_iy+f_i\}.
\]
The first factor is nonnegative on the feasible center interval. The
second is affine on the current interval, so at most one additional
rational root determines the preferred endpoint; an identically zero
factor permits either endpoint. This also gives an affine minimizing
rule after refinement.

Take the union of all leaf breakpoints. Substituting the selected affine
rules into~\eqref{eq:star} gives a rational quadratic on every closed
piece. The rule remains feasible and minimizing at its endpoints:
bound switches coincide there, and endpoint-choice switches have equal
values there. Minimize each quadratic at its endpoints and, for a
positive square coefficient, its feasible stationary point. These
candidates are rational and exhaustive. The selected center and leaf
rules supply an attained rational minimizer. A singleton center interval
uses the same conditional scalar minimization at that point.
\end{proof}

If leaf $i$ has $m_i$ envelope lines, their pair intersections produce
at most $O(m_i^2)$ center breakpoints; a constant number of further
switches per envelope interval gives the same order. The union of
breakpoints has polynomial size. The present implementation deliberately
uses straightforward pair enumeration and scanning. Its arithmetic
cost is polynomial in the supplied row counts and number of leaves;
no claim of an optimal complexity bound is needed for this small-block
oracle. Unlike general elimination on a deeper graph, this construction
does not repeatedly nest piecewise messages along a tree.

The implementation
\pathref{../theory/quadratic_star.py}{theory/quadratic_star.py}
records the complete center partition, affine leaf rules, quadratic
value functions, and attained minima. Unsupported mixed terms or rows
raise an explicit error instead of being dropped. Replay rebuilds the
same exact object from trusted input. As with the polygon oracle, this
is exact support for a supplied direction; completeness of the
integration's direction search is a separate question.

\begin{corollary}[Sharp gain for specified pair directions]\label{cor:gain}
Let compact rational center--leaf polygons $P_i$ assemble into a nonempty
star, and fix rational quadratic pair objectives $q_i(y,x_i)$. Define
\[
\beta_i=\min_{P_i}q_i,\qquad
\beta_* = \min_{(y,x_i)\in P_i\ \forall i}\sum_iq_i(y,x_i),
\qquad \Delta=\beta_*-\sum_i\beta_i.
\]
Then $\Delta\ge0$ is the largest possible constant improvement in
the sum of these fixed pair inequalities. It is exactly computable by
the polygon and star support algorithms. Moreover, $\Delta>0$ if and
only if the center projections of the pair minimizing sets have no
common point.
\end{corollary}
\begin{proof}
Every pair objective is at least its individual minimum, proving
$\Delta\ge0$. Equality holds precisely when one feasible assembled
point minimizes every pair objective. Such a point exists precisely
when their minimizing center sets intersect, since leaves can then
be chosen independently. If no such point exists, compactness makes
the attained assembled minimum strictly larger. The attained value
$\beta_*$ is the largest valid constant in this summed direction.
\end{proof}

This gives a precise diagnostic for a proposed merger of directions.
It does not choose the normals or predict whether computing the gain
will improve total runtime. In the example below, both pair minima
are zero and the exact gain is $1/128$.

The box-only case is a specialization of Del Pia and Khajavirad's
forest-QP algorithm, which gives an $O(n^2)$ arithmetic bound for general
forests \cite{DelPiaKhajavirad2026}. Piecewise-affine responses of
strictly convex parametric quadratic programs are also classical
\cite{BemporadEtAl2002}. The additional domain handled explicitly here
consists of center--leaf affine rows with all signs of scalar leaf
curvature. That scope distinction and its exact implementation are
documented; a first-result claim does not follow from a bounded source
search.
